%! Author = alexpayne
%! Date = 9/4/24

% Preamble
\documentclass[12pt, english, course]{lecture}

% Packages
\usepackage{amsmath}

\title{A Beginners Guide to Scientific Programming in the Chodera Lab}
\author{Alex Payne}
\email{alex.payne@choderalab.org}

% Document
\begin{document}
\section{The Chodera Lab Philosophy of Scientific Programming}
\begin{itemize}
    \item Non-reproducible science is just sparkling philosophy.
    \item Computational science is only reproducible if the 1) data and 2) algorithms are accessible.
    \item Hiding code behind a paywall or making it only "available upon request" is detrimental to scientific achievement.
    \item Results are only useful if they come with confidence intervals.
\end{itemize}

\section{Languages and Computing Environments}
Like most scientific software groups, we mainly use Python, usually as a wrapper for some core C or C++ code that only a few people actually know how to update.
This means we tend to use an IDE such as PyCharm or VSCode, as the automatic detection and annotation of Python syntax is super useful.

Pretty much any code that we write is expected to end up in a Github repo if anyone is expected to ever use it or see it referenced in a talk or a paper.
We do not wait until publication to make our code public, unless it's code for running experiments that contains proprietary information (such as free energy calculations for a collaborator's undisclosed ligand).

\section{The balance between building packages and annotating experiments}
All of the code we write falls into one of these (partially overlapping categories).

\textbf{Library Code}: Classes and functions, organized in modules with a package we are developing

\textbf{Scripting}: Scripts which use library code to perform particular tasks, usually with various arguments enabled by argparse or click.

\textbf{Scientific Code}: Anything which is used to collect and analyze the scientific data needed to run experiments, usually a combination of data files, README's, scripts, and jupyter notebooks.

These three categories, although they blend quite often into each other, have distinct development priorities.
For instance, while library code shouldn't contain filepaths, and likely shouldn't have both code for data collection and code for data analysis

Venn Diagram
    1 - Should be focused; optimized for performance; should be tested with continuous integration on multiple infrastructures to maintain usability.
    2 - optimized for readability; likely won't be tested
    3 - Should contain the data or a way to download the data; optimized for being a concise and accurate description of the requirements to reproduce the experiments; no CI but instead the environment and data should be made available
    1+2 - Shouldn't contain filepaths
    1+3
    2+3 - Can perform a variety of tasks
    1+2+3



\end{document}